% Paper 3 — Methods Section
% Building Local-Language Economic Narrative Indices

\section{Methods}\label{sec:methods}

\subsection{Pipeline Architecture}

The BENI pipeline transforms raw newspaper articles into a validated economic narrative index through five stages:

\begin{enumerate}
    \item \textbf{Data ingestion}: Load BENI v1 article records with publication dates, source identifiers, harmonised categories, and article text.
    \item \textbf{Preprocessing}: Normalize text, remove noise, standardize encoding, hash article text, and flag exact or near-exact duplicates.
    \item \textbf{Classification}: Assign economic relevance scores using weak labels, supervised baselines, and LLM-assisted reference labels.
    \item \textbf{Index construction}: Aggregate monthly scores into the narrative index.
    \item \textbf{Validation}: Correlate with macroeconomic indicators and compare across models.
\end{enumerate}

Each stage is modular and can be replaced or extended without modifying the others. The pipeline is implemented in Python with dependencies limited to standard scientific libraries (pandas, scikit-learn, PyTorch, HuggingFace Transformers).

\subsection{Data}

\subsubsection{BENI v1 Database}\label{sec:data}

The empirical base for the method paper is BENI v1, a harmonised Bangla news database constructed from two upstream sources. The first source is the Potrika Bangla News Corpus \citep{mendeley2024potrika}, which provides dated newspaper articles from 2014--2020. The second source is BNAD, which extends coverage after 2020. The current BENI v1 release, `BENI\_unified\_v1.0\_frozen', contains 1,467,705 harmonised article records: 471,723 from Potrika and 995,982 from BNAD.

Each article record preserves the fields required for reproducible measurement: `article\_id', `dataset\_source', `source\_file', `newspaper', `publication\_date', `year\_month', `category\_original', `category\_harmonised', `headline', `text', `text\_clean', `language', `text\_hash', duplicate flags, economic seed labels, model prediction fields, and release version. This design separates the database layer from the modeling layer: raw article provenance, harmonised metadata, and deduplication fields remain available even when the classifier or index specification changes.

The current merge rule is simple and auditable: Potrika dated 2014--2020 plus BNAD post-2020. The merged file flags 14,195 duplicate rows through text hashing and duplicate grouping. Because the purpose of BENI is monthly economic measurement, dated articles and source-level balance are more important than maximising raw article count. The final index-generation step therefore uses the deduplicated, dated subset of BENI v1 rather than every raw row indiscriminately.

\subsubsection{Calibration Subset}

The existing model-comparison and macro-validation results reported below are legacy calibration results from the dated Potrika subset. Potrika is a publicly available dataset of Bangla-language newspaper articles published between 2014 and 2020. The calibration experiment uses six newspapers:

\begin{itemize}
    \item Jugantor
    \item Ittefaq
    \item Kaler Kontho
    \item Inqilab
    \item Jaijaidin
    \item Somoyer Alo
\end{itemize}

Articles span eight categories: Economy, National, Politics, World News, Sports, Education, Entertainment, and Science \& Technology. For the calibration classification task, Economy articles are treated as positive weak-label examples and National, Politics, and World News articles as negative weak-label examples, down-sampling negatives at a 1:3 ratio to address class imbalance.

The resulting calibration dataset contains 120,707 articles (38,489 Economic, 82,218 non-Economic), split temporally into train (before 2019), validation (2019), and test (2020) sets. These results are retained as prototype evidence. The current frozen BENI v1 validation set contains 186 clean reference labels mapped to canonical BENI v1 article IDs, with 69 Economic and 117 Not Economic examples after excluding the 114 ambiguous rows from clean validation.

\subsubsection{LLM-Generated Silver Standard}

We generated a silver-standard evaluation set of 299 articles from the BNLP News Categorization dataset \citep{banglanlp2021} using Claude Sonnet 4 (claude-sonnet-4-20250514) as an LLM annotator. Articles were selected using stratified sampling to ensure representation across topic categories. The full annotation prompt is provided in Appendix~\ref{app:annotation-prompt}. For the frozen BENI v1 comparison, these labels are mapped to canonical BENI v1 article IDs and filtered to the 186-row clean validation subset.

Each article was annotated using the following protocol:

\begin{itemize}
    \item \textbf{Provider}: Anthropic Claude (claude-sonnet-4-20250514)
    \item \textbf{Temperature}: 0.0 (deterministic)
    \item \textbf{Max tokens}: 500 per response
    \item \textbf{Rate limiting}: 2-second pause every 20 API calls
    \item \textbf{Retry}: Up to 3 attempts on parse failure
\end{itemize}

The annotation schema included six fields, elicited via a structured JSON prompt (reproduced in Appendix~\ref{app:annotation-prompt}):

\begin{enumerate}
    \item \textbf{Economic relevance}: Binary (Economic / Not Economic)
    \item \textbf{Confidence}: Ordinal (1 = guessing, 2 = fairly sure, 3 = certain)
    \item \textbf{Difficulty}: Clear-cut / Borderline (for Economic articles)
    \item \textbf{Economic topic}: 12 categories (inflation, exchange rate, monetary policy, fiscal policy, trade, employment, agriculture, industry, remittances, energy, general, other)
    \item \textbf{Sentiment}: Negative / Neutral / Positive
    \item \textbf{Narrative force}: 8 frames (crisis, burden, blame, reform, stability, uncertainty, resilience, neutral)
\end{enumerate}

The prompt explicitly handled Bangla-specific ambiguities, such as the word \textit{kor}, which can mean ``tax'' in economic contexts or serve as a verb suffix in non-economic contexts. Articles about Indian regional politics (Kolkata, West Bengal) were explicitly instructed to be classified as Not Economic for BENI since the index targets Bangladesh.

Of the 299 annotated articles, 22 (7.4\%) were classified as Economic and 277 (92.6\%) as Not Economic. The annotator reported confidence of 3 (certain) for 286 articles (95.7\%) and 2 (fairly sure) for 13 articles (4.3\%).

\subsubsection{Self-Consistency Evaluation}

Since LLM-generated labels are a proxy for human annotation, we measured the LLM's self-consistency by re-annotating a random subset of 50 articles in a second pass (different seed, otherwise identical protocol). This tests whether the annotator produces stable labels for the same input.

\begin{table}[h]
\centering
\begin{tabular}{lr}
\toprule
Metric & Value \\
\midrule
Observed agreement & 96.0\% \\
Cohen's $\kappa$ & 0.48 \\
Economic (pass 1) & 1 article \\
Economic (pass 2) & 3 articles \\
Disagreements & 2 / 50 \\
\bottomrule
\end{tabular}
\caption{Self-consistency of Claude Sonnet 4 annotations (n=50).}
\end{table}

The 96\% observed agreement indicates high practical consistency. The moderate $\kappa$ (0.48) is partly a prevalence effect \citep{feinstein1990high}: with only 1--3 Economic articles in 50, even 2 disagreements produce a low $\kappa$ despite 48/50 identical labels. Both disagreements involved borderline articles where the economic framing was debatable---one about political party funding (possibly fiscal policy) and one about mobile phone discounts (possibly trade/commerce). This self-consistency ceiling ($\kappa = 0.48$) bounds the maximum agreement any trained model can achieve against this silver standard.

We note that these are LLM-generated labels, not human expert annotations. However, prior work on LLM annotators \citep{gilardi2023chatgpt, hemmingsen2024llm} has shown that frontier models can match or approach human inter-annotator agreement on text classification tasks. We report the self-consistency metric to establish the reliability ceiling of our labels. Future work with human annotators would be needed to establish true gold-standard agreement rates.

\subsubsection{Macroeconomic Indicators}

We validate BENI against three macroeconomic indicators for Bangladesh:

\begin{itemize}
    \item \textbf{BDT/USD exchange rate}: Monthly end-of-period spot rate from the Bank for International Settlements (BIS).
    \item \textbf{Consumer Price Index (CPI)}: Monthly all-items index from the International Monetary Fund (IMF) SDMX API.
    \item \textbf{Foreign exchange reserves}: Annual total reserves from the World Bank API.
\end{itemize}

All series are publicly available without authentication and can be downloaded programmatically.

\subsection{Model 1: TF-IDF + Logistic Regression}

The first model combines TF-IDF vectorization with logistic regression:

\begin{itemize}
    \item \textbf{Vectorizer}: TfidfVectorizer with max\_features = 80,000, min\_df = 2, ngram\_range = (1, 2), sublinear\_tf = True.
    \item \textbf{Classifier}: OneVsRestClassifier wrapping LogisticRegression with class\_weight = ``balanced'', solver = ``liblinear'', max\_iter = 1,000.
\end{itemize}

This model requires no GPU, trains in under two minutes, and serves as a computationally efficient baseline. As we show in Section~\ref{sec:results}, it achieves the highest F1 among weak-label models, despite its architectural simplicity.

\subsection{Model 2: BanglaBERT Trained on Weak Labels}

The second model fine-tunes BanglaBERT \citep{chowdhury2020banglabert} on the same keyword-derived pseudo-labels as Model~1. BanglaBERT is an Electra-based pre-trained language model for Bangla (424M parameters), available on HuggingFace. This model is retained as legacy prototype evidence; the frozen BENI v1 validation comparison reported in this paper uses TF-IDF + Logistic Regression as the production baseline.

\subsubsection{Weak Label Generation}

We generate pseudo-labels using a keyword matching approach. A list of 32 Bangla economic terms (Appendix~\ref{app:keywords}) is matched against the normalized text of each article. Articles containing at least one keyword are labeled as Economic; others as Not Economic.

On the calibration subset, this produces a training set of 120,707 articles with approximately 4.4\% positive rate---noisy but sufficient for initial fine-tuning and for diagnosing the weak-label ceiling. That calibration result is retained as a prototype appendix, not as the final BENI v1 validation result.

\subsubsection{Fine-Tuning Protocol}

\begin{itemize}
    \item \textbf{Optimizer}: AdamW with learning rate = 2e-5.
    \item \textbf{Schedule}: Linear warmup (10\% of steps) followed by linear decay.
    \item \textbf{Batch size}: 32 (Kaggle T4 GPU).
    \item \textbf{Max sequence length}: 384 tokens.
    \item \textbf{Epochs}: 3.
    \item \textbf{Mixed precision}: fp16 with gradient clipping (max\_norm = 1.0).
\end{itemize}

Training requires approximately 2--3 hours on a Kaggle T4 GPU (16 GB VRAM). The best model is selected by validation macro F1.

\subsection{Model 3: LLM-Generated Silver Standard}

The third approach uses the LLM-generated silver standard directly as the gold standard for evaluation. Rather than treating the 299 LLM-annotated articles as training data for a downstream classifier, we use the LLM annotations themselves as the reference labels against which Models~1 and~2 are evaluated. For the frozen BENI v1 paper, the relevant clean validation subset is the 186-row version of this silver standard, not the full 299-row prototype set.

\subsection{Active Learning Simulation}\label{sec:active-learning}

We measure how classification performance scales with annotation investment using active learning simulation. Because the silver standard contains only 22 positive examples (7.4\% of 299), we use TF-IDF + Logistic Regression as the classifier---the best-performing weak-label model from the prototype calibration results---rather than BanglaBERT, which cannot learn the minority class from this few examples. The protocol uses 5-fold stratified cross-validation to maximize the effective test size:

\begin{enumerate}
    \item Split the 299 LLM-annotated articles into 5 stratified folds preserving the 7.4\% economic proportion.
    \item For each annotation budget $k \in \{50, 100, 150, 200, 250, 299\}$:
    \begin{enumerate}
        \item For each fold, train the classifier on the training articles of that fold, subsampled to $k$ articles (stratified by economic relevance).
        \item Evaluate on the held-out fold.
        \item Repeat 5 times per fold with different random seeds, yielding 25 estimates per budget.
    \end{enumerate}
    \item Report mean and standard deviation of F1 (Economic class) and accuracy across all trials.
\end{enumerate}

This analysis answers the practical question: \textit{how many annotated articles are needed for reliable economic text classification, and does TF-IDF overcome the base-rate problem at any feasible budget?}

\subsection{Narrative Index Construction}

The BENI v1 monthly index is constructed as follows:

\begin{enumerate}
    \item Run the trained classifier on the deduplicated, dated BENI v1 article records to obtain economic probability scores.
    \item Assign binary predictions: $\text{economic\_pred}_i = \mathbb{1}[\text{prob}_i > 0.5]$.
    \item Aggregate by month and source: compute economic share (the proportion of articles predicted as Economic) and mean probability within each source-month.
    \item Produce both article-weighted and source-balanced monthly indices. The source-balanced version is preferred for inference because it limits domination by high-volume outlets.
    \item Preserve model version, release version, and index specification so that later validation and nowcasting papers can use a frozen BENI series.
\end{enumerate}

The prototype correlations reported in Section~\ref{sec:results} use the older 120,707-article Potrika calibration index. They should be interpreted as validation diagnostics for the pipeline, not as the final BENI v1 macro-validation results. The frozen monthly index outputs generated from the repaired BENI v1 article predictions are the ones used for the current release package.

\subsection{Validation}

We validate BENI using three approaches:

\begin{enumerate}
    \item \textbf{Model comparison}: Compare accuracy, macro F1, and weighted F1 on the frozen 186-row BENI v1 validation set. Use the repaired TF-IDF predictions as the production baseline; retain the BanglaBERT comparison only as prototype evidence.
    \item \textbf{Active learning curve}: Report F1 (Economic class) and accuracy vs. annotation budget with standard deviations from 5-fold cross-validation on the legacy prototype set.
    \item \textbf{Macro correlation}: Compute Pearson and Spearman correlations between the monthly BENI index and macroeconomic indicators (FX, CPI, reserves) on both raw levels and first-differenced series.
\end{enumerate}

The first-differenced analysis is critical: level correlations can be spurious (both series trend over time), while first-differenced correlations test whether month-to-month changes in economic news share predict month-to-month changes in macroeconomic conditions.
